\begin{abstract}
AI systems can now propose ideas, modify code, run experiments, and write research reports. The harder question is whether they make better research decisions: choosing useful unknowns, designing informative tests, rejecting weak explanations, and carrying lessons into new projects. This survey organizes 79 resources from the Awesome Autoresearch collection around five questions: research taste, question formulation, reliable execution, learning from experience, and evaluation. It compares the feedback and evidence used by different methods rather than treating workflow completion as a single measure of scientific ability. The collection contains 72 papers, six essays or talks, and one software project. Each question section ends with limitations, open questions, and two concrete experimental directions. We then identify three practical starting projects and a shared evaluation protocol. AI research is the main testbed, with mathematics as a complementary setting. The proposed experiments are research designs, not reported results or claims of established novelty.
\end{abstract}

\noindent\textbf{Keywords:} autoresearch; AI for meta-science; research taste; hypothesis generation; AI scientists; research-agent learning.
